\documentclass[10pt,a4paper]{amsart}
\usepackage[margin=19mm]{geometry}
\usepackage{amsmath,amssymb,mathtools}
\usepackage[scr=rsfs,cal=euler]{mathalfa}
\usepackage{xfrac}
\usepackage[unicode,hidelinks,bookmarks=false]{hyperref}
\newcommand{\ie}{i.e.\ }
\newcommand{\cf}{cf.\ }
\newcommand{\resp}{respectively\ }
\newcommand{\Subsection}{\subsection}
\providecommand{\nobreakdash}{\nobreak\hbox{-}}
\setlength{\emergencystretch}{2em}
\multlinegap0pt
\usepackage{bm}
\usepackage{bbm}
\usepackage{dsfont}
\usepackage{xspace}
\usepackage{cancel}
\usepackage{mathtools}
\usepackage{amsthm}
\makeatletter
\@ifundefined{theorem}{\newtheorem{theorem}{Theorem}}{}
\@ifundefined{lemma}{\newtheorem{lemma}[theorem]{Lemma}}{}
\makeatother
\def\eqref#1{equation~\ref{#1}}
\newcommand{\disc}{\mathrm{Disc}}
\newcommand{\dtv}{\mathrm{d_{TV}}}
\title[Testing with Non-identically Distributed Samples]{Testing with Non-identically Distributed Samples}
\author{Shivam Garg, Chirag Pabbaraju, Kirankumar Shiragur, Gregory Valiant}
\date{Source: arXiv:2311.11194v2 (CC BY 4.0)}
\begin{document}
\maketitle

\noindent\emph{Source and licence.} Testing with Non-identically Distributed Samples. Version arXiv:2311.11194v2 (CC BY 4.0), distributed under the Creative Commons Attribution 4.0 International licence, as recorded in the arXiv licence metadata for this exact version and shown on the abstract page https://arxiv.org/abs/2311.11194v2. Full rights notice: licenses/arxiv-2311.11194-CC-BY-4.0.txt.\par\smallskip

\noindent\textit{Editorial scope.} Counted unit: the Lemma whose source label is \texttt{lem:disc-conv-commute-gauss} in the article (source0.tex lines 1826--1835, source label \texttt{lem:disc-conv-commute-gauss}), delivered together with the article's own complete original proof of it (source0.tex lines 1837--1910, one \texttt{proof} environment of 74 lines). No mathematics is written, repaired or re-derived: statement and proof are the article's own text line for line. Required context retained uncounted: lem:shift-L1 (lines 1913--1920). Every definition, display and label of those lines is retained unchanged. Omitted from this package and not redistributed: 1--1825, 1836--1836, 1911--1912, 1921--2531. Nothing inside a retained excerpt is omitted or altered: every retained body is delivered exactly as the article's lines are written, so no omission receipt of any kind accompanies this package. Citations: the retained excerpts carry no citation command at all, so no bibliography entry of the article is transcribed and no reference is redistributed. Reference adaptation: none was needed and none was made; every reference inside the delivered unit resolves inside it, so no unresolved reference or citation is produced. Printed slips kept literal: no printed word, symbol, hypothesis or number of the counted statement or of its proof was altered. Layout and class: the article's own class is \texttt{article} with packages including tmlr, amsmath, amsfonts, bm, hyperref, url, microtype, graphicx, subcaption, caption, booktabs, inputenc, fontenc, nicefrac; the wrapper is the lane's pinned \texttt{amsart} editorial wrapper at 10pt on A4, which restates the theorem-like environments the retained text needs and adds the article's own macro definitions that the retained text needs (\texttt{\textbackslash disc}, \texttt{\textbackslash dtv}, \texttt{\textbackslash eqref}), with extra wrapper packages (bm, bbm, dsfont, xspace, cancel, mathtools). The delivered numbering is the unit's own. Assets: no retained span contains a figure environment, a graphics-inclusion command, a PDF-page inclusion, a table or a TikZ picture; no image file is copied into this package, the package declares no asset dependency and no image file is redistributed with it. Licence: version arXiv:2311.11194v2 of this article is distributed under the Creative Commons Attribution 4.0 International licence, as recorded in the arXiv licence metadata for this exact version and shown on the abstract page https://arxiv.org/abs/2311.11194v2. A full rights notice is delivered at licenses/arxiv-2311.11194-CC-BY-4.0.txt. Characters: every retained excerpt is pure ASCII, so the delivered engine, XeLaTeX with its default Latin Modern text font, reports no missing character.
\begin{lemma}
\label{lem:shift-L1}
Let $h:\mathbb{R}^d\to\mathbb{R}_+$ be locally absolutely continuous with $\nabla h\in L^1(\mathbb{R}^d)$.
Then for every $u\in\mathbb{R}^d$,
\[
\|h(\cdot+u)-h(\cdot)\|_{L^1}\;\le\;\|u\|_2\,\|\nabla h\|_{L^1}.
\]
\end{lemma}

\begin{lemma}[Discretization--convolution almost commute for Gaussians]
\label{lem:disc-conv-commute-gauss}
Let $X\sim\mathcal{N}(\mu_X,\Sigma_X)$ and $Y\sim\mathcal{N}(\mu_Y,\Sigma_Y)$ be independent in $\mathbb{R}^d$ (with $d$ constant), and assume at least one of $\Sigma_X,\Sigma_Y$ is positive definite.\footnote{If one covariance is singular, the bound below with that matrix should be read as $+\infty$; the inequality is then applied with the other (positive definite) covariance.} Let $\disc(\cdot)$ be coordinate-wise rounding to the nearest integer (ties are immaterial since Gaussians put zero mass on boundaries). Then
\[
\dtv\!\big(\disc(X)+\disc(Y),\,\disc(X+Y)\big)
\;\le\; \frac{\sqrt{d}}{4}\,
\min\!\Big\{\sqrt{\mathrm{Tr}(\Sigma_X^{-1})},\ \sqrt{\mathrm{Tr}(\Sigma_Y^{-1})}\Big\}.
\]
In particular, if $\lambda_{\min}(\Sigma_X)\to\infty$ or $\lambda_{\min}(\Sigma_Y)\to\infty$, the left-hand side is $O(1/\sqrt{\lambda_{\min}})$ and tends to $0$.
\end{lemma}

\begin{proof}
\emph{Without loss of generality assume $\Sigma_Y\succ 0$; otherwise swap $X$ and $Y$.}
Let $f$ and $g$ be the densities of $X$ and $Y$. For $m\in\mathbb{Z}^d$, write $C_m:=m+[-\tfrac12,\tfrac12)^d$.
Define
\[
P(m):=\Pr(\disc(X)+\disc(Y)=m)
=\sum_{a\in\mathbb{Z}^d}\Big(\int_{C_a} f(x)\,dx\Big)\Big(\int_{C_{m-a}} g(y)\,dy\Big),
\]
\[
Q(m):=\Pr(\disc(X+Y)=m)
=\int_{C_m}\int_{\mathbb{R}^d} f(x)\,g(z-x)\,dx\,dz.
\]
For $x\in C_a$, set $\delta(x):=a-x\in[-\tfrac12,\tfrac12]^d$. Changing variables $z=y+a$ in $Q(m)$ yields
\[
Q(m)=\sum_{a\in\mathbb{Z}^d}\int_{C_a}\!\Big(\int_{C_{m-a}} g\big(y+\delta(x)\big)\,dy\Big) f(x)\,dx,
\]
and hence
\begin{equation}\label{eq:PmQm-diff}
P(m)-Q(m)
=\sum_{a\in\mathbb{Z}^d}\int_{C_a}\!\Big(\int_{C_{m-a}}\!\big[g(y)-g\big(y+\delta(x)\big)\big]\,dy\Big) f(x)\,dx.
\end{equation}

\paragraph{Bounding $\sum_m |P(m)-Q(m)|$ by an expectation times a norm.}
Summing \eqref{eq:PmQm-diff} in absolute value over $m$ and using that the cubes $\{C_{m-a}\}_m$ are disjoint and partition $\mathbb{R}^d$,
\begin{align}
\sum_{m\in\mathbb{Z}^d}\!|P(m)-Q(m)|
&\le \sum_{a}\int_{C_a} f(x)\,\Big(\int_{\mathbb{R}^d}\!\big|g(y)-g(y+\delta(x))\big|\,dy\Big)\,dx \nonumber\\
&= \int_{\mathbb{R}^d} f(x)\,\big\|g(\cdot)-g(\cdot+\delta(x))\big\|_{L^1}\,dx.
\label{eq:sum-bound-1}
\end{align}
(\emph{We used }$\sum_m\big|\int_{C_{m-a}}\!\cdot\big|
\le \sum_m\int_{C_{m-a}}\!|\cdot|
= \int_{\mathbb{R}^d}\!|\cdot|$,
since $\{C_{m-a}\}_m$ partition $\mathbb{R}^d$.)

Apply Lemma~\ref{lem:shift-L1} with $h=g$ and $u=\delta(x)$ pointwise in $x$ to get
\begin{equation}\label{eq:sum-bound-2}
\big\|g(\cdot)-g(\cdot+\delta(x))\big\|_{L^1}\ \le\ \|\delta(x)\|_2\ \|\nabla g\|_{L^1}.
\end{equation}
Plugging \eqref{eq:sum-bound-2} into \eqref{eq:sum-bound-1} and factoring out the constant $\|\nabla g\|_{L^1}$ yields
\begin{equation}\label{eq:product-form}
\sum_{m\in\mathbb{Z}^d}\!|P(m)-Q(m)|
\ \le\ \Big(\,\mathbb{E}\big[\|\delta(X)\|_2\big]\,\Big)\ \cdot\ \|\nabla g\|_{L^1}.
\end{equation}
Since $X\in C_{\disc(X)}$ almost surely, $\delta(X)=\disc(X)-X\in[-\tfrac12,\tfrac12]^d$, and thus
\begin{equation}\label{eq:delta-expect}
\mathbb{E}\big[\|\delta(X)\|_2\big]\ \le\ \sup_{\|\delta\|_\infty\le 1/2}\|\delta\|_2\ =\ \frac{\sqrt{d}}{2}.
\end{equation}

\paragraph{Computing $\|\nabla g\|_{L^1}$.}
For $Y\sim\mathcal{N}(\mu_Y,\Sigma_Y)$,
\[
\nabla g(y)=-\,\Sigma_Y^{-1}(y-\mu_Y)\,g(y),
\qquad
\|\nabla g\|_{L^1}=\int \|\Sigma_Y^{-1}(y-\mu_Y)\|_2\,g(y)\,dy
=\mathbb{E}\big[\|\Sigma_Y^{-1/2}Z\|_2\big],
\]
where $Z\sim\mathcal{N}(0,I_d)$. By Cauchy--Schwarz,
\begin{equation}\label{eq:grad-L1}
\|\nabla g\|_{L^1}\ \le\ \sqrt{\mathbb{E}\!\left[Z^\top \Sigma_Y^{-1} Z\right]}\ =\ \sqrt{\mathrm{Tr}(\Sigma_Y^{-1})}.
\end{equation}

Combining \eqref{eq:product-form}, \eqref{eq:delta-expect}, and \eqref{eq:grad-L1},
\[
\sum_{m\in\mathbb{Z}^d}\!|P(m)-Q(m)|\ \le\ \frac{\sqrt{d}}{2}\,\sqrt{\mathrm{Tr}(\Sigma_Y^{-1})},
\]
hence
\[
\dtv\!\big(\disc(X)+\disc(Y),\,\disc(X+Y)\big)
=\tfrac12\sum_{m}|P(m)-Q(m)|
\ \le\ \frac{\sqrt{d}}{4}\,\sqrt{\mathrm{Tr}(\Sigma_Y^{-1})}.
\]
By symmetry, the same bound holds with $\Sigma_X$; taking the minimum proves the claim.
\end{proof}

\end{document}
